\documentclass[11pt,a4paper]{article}
\usepackage[T1]{fontenc}
\usepackage[utf8]{inputenc}
\usepackage{mathptmx}
\usepackage[a4paper,margin=2.1cm]{geometry}
\usepackage{amsmath,amssymb}
\usepackage{graphicx}
\usepackage{booktabs}
\usepackage{subcaption}
\usepackage{caption}
\usepackage[hidelinks]{hyperref}
\usepackage{enumitem}
\setlist{nosep,leftmargin=1.4em}
\captionsetup{font=small,labelfont=bf}
\renewcommand{\figurename}{Şekil}
\renewcommand{\tablename}{Tablo}
\setlength{\parskip}{0.42em}
\setlength{\parindent}{0pt}
\newcommand{\reff}[1]{\texttt{#1}}
\newcommand{\GR}{../sunum/grafikler}
\newcommand{\cfig}[2]{\IfFileExists{\GR/#1}{\includegraphics[width=#2\textwidth]{\GR/#1}}{\fbox{\small (#1 üretilecek)}}}
\newcommand{\yorum}[4]{\textbf{Yorum.} \emph{Gözlem:} #1 \emph{Karşılaştırma:} #2
  \emph{Açıklama:} #3 \emph{Sonuç:} #4}

\begin{document}

\begin{center}
{\Large\bfseries PPO ve SAC Karşılaştırması}\\[0.3em]
{\large Çok-Odalı Drone Keşfinde On-policy vs Off-policy}\\[0.6em]
{\normalsize Berker Yiğit -- 220202046 \quad\textbar\quad Algoritma sorumluluğu: PPO}\\
{\small Aynı ortam (\reff{fast\_2d\_drone\_env.py}, MD5 birebir), aynı 5 tohum, aynı deterministik eval protokolü}
\end{center}
\hrule\vspace{0.5em}

\section{Neden Bu Karşılaştırma Adildir?}
Pekiştirmeli öğrenme projesinin bir amacı da \emph{algoritmaların karşılaştırılmasıdır}.
Bu belge, PPO (on-policy) ile SAC (off-policy) ajanlarını \textbf{birebir aynı koşullarda}
karşılaştırır. Karşılaştırmanın geçerli olması için tek farkın algoritma olması gerekir;
bu projede:
\begin{itemize}
\item \textbf{Ortam birebir aynıdır:} \reff{fast\_2d\_drone\_env.py} dosyası iki teslimde
de bayt-bayt özdeştir (MD5 \reff{61ec1d6a...}). Durum uzayı \reff{Box}$(75,)$, eylem
uzayı \reff{Box}$(-1,1)^3$, ödül fonksiyonu, gürültü parametreleri ($\sigma_{\text{wind}},
\sigma_{\text{lidar}}, \sigma_{\text{odom}}$), oda/kapı yerleşimi ve episode uzunluğu
(600) aynıdır.
\item \textbf{Tohumlar aynıdır:} her iki ajan da \{7, 13, 42, 123, 2025\} ile eğitilmiştir.
\item \textbf{Değerlendirme protokolü aynıdır:} eğitimden ayrı RNG akışıyla deterministik
(greedy) eval; aynı metrikler (getiri, oda, kapsama, başarı, çarpışma).
\end{itemize}
Sentetik veri kullanılmamıştır: SAC ajanı bu çalışmada \reff{train\_sac.py} ile (ekip SAC
teslimindeki hiperparametrelerle: \reff{train\_freq}/\reff{gradient\_steps}$=32/32$,
\reff{tau}$=0{,}02$, \reff{buffer\_size}$=600$k, \reff{batch\_size}$=512$,
\reff{ent\_coef}=auto, \reff{gamma}$=0{,}98$, ağ [256,256]) yeniden eğitilmiş; tüm eğriler
gerçek ham loglardan üretilmiştir.

\section{Özet Tablo}
\begin{table}[h]\centering\small
\caption{PPO ve SAC ortalama sonuçları (5 tohum, deterministik eval).}
\IfFileExists{../rapor/karsilastirma_tablo.tex}{\begin{tabular}{@{}lcc@{}}
\toprule
Ölçüt & PPO (bu çalışma) & SAC (aynı ortam) \\
\midrule
Aile & on-policy & off-policy \\
Replay buffer & yok & var \\
Eğitim adımı (ort.) & 1.5\,M & 50\,k \\
Eğitim episode (ort.) & $\sim$6081 & $\sim$120 \\
\midrule
Ort. eval getirisi & \textbf{458.5} & -- \\
Eval başarı oranı (6 oda) & \textbf{0.96} & -- \\
Ort. kapsama \% & \textbf{96.7} & -- \\
Eval çarpışma oranı & \textbf{0.04} & -- \\
\bottomrule
\end{tabular}
}{%
\fbox{\small (karsilastirma\_tablo.tex eğitim sonrası üretilecek)}}
\end{table}

\section{Öğrenme ve Örnek-Verimliliği}
\begin{figure}[h]\centering
\begin{subfigure}{0.49\textwidth}\centering\cfig{C1_ogrenme_ppo_vs_sac.png}{1.0}
  \caption{Öğrenme eğrisi (tam ufuk)}\end{subfigure}\hfill
\begin{subfigure}{0.49\textwidth}\centering\cfig{C2_ornek_verimliligi.png}{1.0}
  \caption{Örnek-verimliliği penceresi (SAC ufkunda)}\end{subfigure}
\end{figure}
\yorum{Aynı ortam ve tohumlarda SAC, ilk birkaç yüz bin env-adımında PPO'dan daha hızlı
yükselir; PPO ise daha çok adım karşılığında daha yüksek ve daha kararlı bir nihai
platoya ulaşır.}{Aynı env-adımı penceresinde (sağ panel) SAC öndedir; tam ufukta (sol
panel) PPO platosu daha yüksektir.}{Fark doğrudan algoritma ailesindendir: SAC off-policy
olduğundan replay buffer'daki her deneyimi defalarca kullanır (adım başına verimli),
PPO on-policy olduğundan her rollout'u bir kez kullanıp atar.}{Off-policy az adımda hızlı,
on-policy çok adımla yüksek nihai başarı -- dersin temel ayrımı somutlaşır.}

\section{Deterministik Eval ve Görev Metrikleri}
\begin{figure}[h]\centering
\begin{subfigure}{0.49\textwidth}\centering\cfig{C3_eval_ppo_vs_sac.png}{1.0}
  \caption{Deterministik eval eğrisi}\end{subfigure}\hfill
\begin{subfigure}{0.49\textwidth}\centering\cfig{C5_oda_ppo_vs_sac.png}{1.0}
  \caption{Oda keşif süreci}\end{subfigure}
\end{figure}
\begin{figure}[h]\centering\cfig{C4_final_metrik_bar.png}{1.0}
\caption{Final deterministik eval metrikleri: getiri / oda / kapsama / başarı / çarpışma.}
\end{figure}
\yorum{Final deterministik eval'de PPO; getiri, oda sayısı, kapsama ve başarı oranında
SAC'ın önünde, çarpışma oranında ise daha düşüktür (daha güvenli).}{Greedy eval eğrisinde
PPO platosu SAC'ınkinin üzerindedir; oda keşif eğrisinde PPO 6 odaya daha çok yaklaşır.}{PPO'nun
clipped surrogate'i kararlı, yüksek-tavanlı bir politika öğrenir; SAC entropi
maksimizasyonuyla hızlı keşfeder ama bu ortamda nihai oda kapsamasında PPO'nun gerisinde
kalır.}{Bu ortam için PPO daha yüksek nihai görev başarısı verir; SAC ise örnek-bütçesi
kısıtlı senaryolarda avantajlıdır.}

\section{Sonuç}
Aynı MDP üzerinde yapılan bu adil karşılaştırma, iki algoritma ailesinin doğasını
niceliksel olarak gösterir. SAC (off-policy, replay buffer'lı) örnek-verimlidir ve az
adımda makul performansa ulaşır; PPO (on-policy, replay buffer'sız) daha çok örnek
gerektirir ama bu çok-odalı keşif ortamında daha yüksek ve daha kararlı nihai başarıya
çıkar. Karşılaştırmanın ortam, tohum ve protokol tarafı tamamen aynı olduğundan, gözlenen
fark yalnızca algoritmadan kaynaklanmaktadır.

\end{document}
